% ============================================================
%  CONCLUSION GÉNÉRALE
% ============================================================
\chapter*{Conclusion générale}
\addcontentsline{toc}{chapter}{Conclusion générale}
\markboth{Conclusion générale}{Conclusion générale}

Ce projet de fin d'études m'a permis de concevoir et de développer une plateforme web interactive complète dédiée à la valorisation de l'huile d'olive extra vierge du Domaine FENDRI. Depuis la modélisation 3D des contenants sous Blender jusqu'au déploiement d'une application e-commerce multilingue et sécurisée, chaque étape de ce travail a représenté un défi technique que j'ai dû analyser, comprendre et résoudre.

L'une des expériences les plus formatrices de ce projet a été la confrontation avec les limites réelles des technologies. Face aux difficultés d'intégration des modèles 3D via WebGL, j'ai dû revoir mes choix initiaux, explorer des alternatives et finalement adopter une solution basée sur des rendus WebP haute définition, qui s'est révélée plus performante et plus fidèle visuellement. Cette démarche m'a appris que la maîtrise d'un projet ne réside pas uniquement dans la mise en œuvre de ce qui a été planifié, mais aussi dans la capacité à s'adapter et à trouver des solutions pertinentes face à l'inattendu.

Sur le plan des résultats, j'ai développé une application structurée autour d'une architecture React et Express.js, intégrant un configurateur interactif en six étapes, un espace client et un espace administrateur complets, un chatbot trilingue avec support RTL, ainsi qu'un système de paiement multi-méthodes. L'ensemble des assets graphiques a été optimisé avec une réduction de poids de 67 à 70\,\%, et l'application a été validée à travers dix-huit tests unitaires sans erreur ESLint ni TypeScript.

Ce stage m'a également permis de travailler dans un cadre professionnel réel au sein de Digilab Solutions, ce qui m'a confrontée aux exigences concrètes d'un projet client : respect des délais, cohérence de l'identité visuelle, qualité du code et soin apporté à l'expérience utilisateur.

En guise de perspectives d'amélioration, plusieurs évolutions pourraient enrichir la plateforme dans sa version future :
\begin{itemize}
  \item L'intégration d'un moteur de recommandation personnalisé basé sur l'historique des achats.
  \item Le déploiement d'un rendu 3D temps réel via Three.js ou Babylon.js, une fois les problèmes de compatibilité des matériaux résolus.
  \item L'extension vers une application mobile native afin de toucher un public plus large.
\end{itemize}

En définitive, ce projet illustre concrètement la manière dont les technologies du numérique peuvent servir la valorisation d'un patrimoine agricole et culturel, en transformant un produit du terroir en une expérience premium personnalisée, à la hauteur des standards internationaux du commerce digital.
